
\chapter{بین‌المللی‌سازی، RTL/LTR و Accessibility}

\section{زبان‌ها}

\begin{decisionbox}
فارسی زبان اصلی و انگلیسی زبان دوم است. هر دو از ابتدای Frontend در معماری حضور دارند.
\end{decisionbox}

\section{UI Language vs Content Language}

دو مفهوم مستقل:
\begin{itemize}
  \item \eng{preferred\_ui\_language}
  \item content language preferences.
\end{itemize}

مثال:
\begin{itemize}
  \item UI = فارسی؛
  \item content = فارسی + انگلیسی.
\end{itemize}

\section{Translation Files}

\begin{latin}
\begin{lstlisting}
src/i18n/
├── fa/
│   ├── common.json
│   ├── auth.json
│   ├── feed.json
│   └── profile.json
└── en/
    ├── common.json
    ├── auth.json
    ├── feed.json
    └── profile.json
\end{lstlisting}
\end{latin}

\section{Hard-coded text}

بعد از راه‌اندازی i18n:
\begin{latin}
\begin{lstlisting}
<button>{t("idea.save")}</button>
\end{lstlisting}
\end{latin}

بهتر از:
\begin{latin}
\begin{lstlisting}
<button>Save</button>
\end{lstlisting}
\end{latin}

\section{Direction}

در Root document:
\begin{latin}
\begin{lstlisting}
<html lang="fa" dir="rtl">
\end{lstlisting}
\end{latin}

برای English:
\begin{latin}
\begin{lstlisting}
<html lang="en" dir="ltr">
\end{lstlisting}
\end{latin}

\section{CSS Logical Properties}

ترجیح:
\begin{latin}
\begin{lstlisting}
margin-inline-start: 1rem;
padding-inline-end: 0.5rem;
text-align: start;
border-inline-start: 2px solid;
\end{lstlisting}
\end{latin}

به جای left/rightهای زیاد.

\section{Persian Text Normalization}

Search فارسی باید به تفاوت‌های:
\begin{itemize}
  \item ی/ي؛
  \item ک/ك؛
  \item نیم‌فاصله؛
  \item فاصله‌های اضافی؛
  \item punctuation
\end{itemize}
توجه کند.

\begin{warningbox}
Original content را برای Search به‌صورت destructive تغییر نده. Normalization را برای Query/indexed representation انجام بده.
\end{warningbox}

\section{Accessibility Target}

منبع WCAG 2.1 AA را می‌خواهد.

Project:
\begin{itemize}
  \item WCAG 2.1 AA minimum target؛
  \item استفاده از Improvementهای WCAG 2.2 در صورت ساده بودن؛
  \item عدم ادعای Compliance کامل بدون Audit واقعی.
\end{itemize}

\section{Semantic HTML}

قاعده:
\begin{itemize}
  \item Button واقعی برای Action؛
  \item Link واقعی برای Navigation؛
  \item Form label؛
  \item heading order؛
  \item main/nav/section/article؛
  \item table برای data table.
\end{itemize}

\section{Keyboard}

تمام Interactionهای مهم:
\begin{itemize}
  \item Tab reachable؛
  \item focus visible؛
  \item Enter/Space behavior استاندارد؛
  \item modal focus trapping اگر modal وجود دارد؛
  \item no hover-only action.
\end{itemize}

\section{Screen Reader}

Elementها باید accessible name داشته باشند.

Icon-only button:
\begin{latin}
\begin{lstlisting}
<button aria-label={t("idea.save")}>
  ...
</button>
\end{lstlisting}
\end{latin}

ARIA فقط وقتی Semantic HTML کافی نیست.

\section{Color/Contrast}

Design prototype dark theme دارد. باید Contrast واقعی:
\begin{itemize}
  \item متن اصلی؛
  \item متن secondary؛
  \item button؛
  \item focus؛
  \item error
\end{itemize}
تست شود.

\section{Touch Targets}

روی موبایل Buttonهای کوچک icon-only باید target لمس کافی داشته باشند، حتی اگر icon visual کوچک باشد.

\section{TTS مستقل از Accessibility}

TTS Feature مصرف محتواست، جای Screen Reader و Accessibility semantic را نمی‌گیرد.

\section{Testing}

\begin{itemize}
  \item keyboard manual؛
  \item browser accessibility tree؛
  \item automated axe-style checks later؛
  \item Playwright critical flows؛
  \item RTL visual inspection.
\end{itemize}

\section{Bilingual Definition of Done}

هر Screen:
\begin{itemize}
  \item fa translation؛
  \item en translation؛
  \item RTL/LTR؛
  \item long text overflow؛
  \item number/date format اگر مرتبط؛
  \item accessible labels.
\end{itemize}
